\frfilename{mt235.tex}
\versiondate{30.3.03/20.8.08}
\copyrightdate{1994}

\def\chaptername{Teorema Radon-Nikod\'ym}
\def\sectionname{Transformasi terukur}
\def\tildeTau{\tilde{\text{T}}}

\newsection{235}

Sekarang saya beralih ke suatu topik yang terpisah dari Teorema
Radon-Nikod\'ym, tetapi tampaknya lebih cocok ditempatkan di sini daripada
di salah satu dari dua bab berikutnya. Saya berusaha memberikan
hasil-hasil yang memperumum rumus dasar kalkulus

\Centerline{$\int g(y)dy=\int g(\phi(x))\phi'(x)dx$}

\noindent dalam konteks suatu transformasi umum $\phi$ di antara
ruang-ruang ukur. Hasil-hasil utamanya, menurut saya, adalah 235A/235E,
yang merupakan dua ungkapan sangat mirip dari gagasan dasar tersebut,
dan 235J, yang memberikan kriteria umum untuk hasil yang lebih kuat.
Suatu perumusan dari arah yang berbeda terdapat dalam 235R.

\leader{235A}{}\cmmnt{ Saya mulai dengan hasil dasar, yang sudah memadai
untuk sebagian besar penerapan yang saya bayangkan.

\medskip

\noindent}{\bf Teorema} Misalkan $(X,\Sigma,\mu)$ dan $(Y,\Tau,\nu)$
ruang-ruang ukur, dan $\phi:D_{\phi}\to Y$,
$J:D_J\to\coint{0,\infty}$ fungsi-fungsi yang didefinisikan pada
subhimpunan koterabaikan $D_{\phi}$ dan $D_J$ dari $X$ sedemikian sehingga

\Centerline{$\int J\times\chi(\phi^{-1}[F])d\mu$ ada $=\nu F$}

\noindent setiap kali $F\in\Tau$ dan $\nu F<\infty$. Maka

\Centerline{$\int_{\phi^{-1}[H]}J\times g\phi\,d\mu$ ada =
$\int_Hg\,d\nu$}

\noindent untuk setiap fungsi yang terintegralkan terhadap $\nu$ dan
dinotasikan dengan $g$, yang bernilai dalam $[-\infty,\infty]$, serta
setiap $H\in\Tau$, dengan ketentuan
bahwa kita
menafsirkan $(J\times g\phi)(x)$ sebagai $0$ ketika $J(x)=0$ dan
$g(\phi(x))$ tidak terdefinisi. Akibatnya, dengan menafsirkan
$J\times f\phi$ dengan cara yang sama,

\Centerline{$\underline{\intop}fd\nu
\le\underline{\intop}J\times f\phi\,d\mu
\le\overline{\intop}J\times f\phi\,d\mu\le\overline{\intop}fd\nu$}

\noindent untuk setiap fungsi bernilai $[-\infty,\infty]$ yang dinotasikan
dengan $f$ dan didefinisikan hampir di mana-mana dalam $Y$.

\proof{{\bf (a)} Jika $g$ suatu fungsi sederhana, katakanlah
$g=\sum_{i=0}^na_i\chi F_i$ dengan $\nu F_i<\infty$ untuk setiap $i$,
maka

\Centerline{$\int J\times g\phi\,d\mu
=\sum_{i=0}^na_i\int J\times\chi(\phi^{-1}[F_i])\,d\mu
=\sum_{i=0}^na_i\nu F_i=\int g\,d\nu$.}

\medskip

{\bf (b)} Jika $\nu F=0$, maka
$\int J\times\chi(\phi^{-1}[F])=0$, sehingga $J=0$ hampir di mana-mana\ pada
$\phi^{-1}[F]$. Jadi, jika $g$ terdefinisi
$\nu$-hampir di mana-mana, maka $J=0\,\,\mu$-hampir di mana-mana\ pada
$X\setminus\dom(g\phi)=(X\setminus D_{\phi})\cup\phi^{-1}[Y\setminus\dom
g]$, dan menurut konvensi yang diusulkan, $J\times g\phi$ terdefinisi
$\mu$-hampir di mana-mana. Selain itu, jika
$\lim_{n\to\infty}g_n=g\,\,\nu$-hampir di mana-mana, maka
$\lim_{n\to\infty}J\times g_n\phi=J\times
g\phi\,\,\mu$-hampir di mana-mana. Jadi, jika $\sequencen{g_n}$ suatu
barisan tak menurun dari fungsi-fungsi sederhana yang konvergen hampir
di mana-mana ke $g$, maka $\sequencen{J\times g_n\phi}$ merupakan
barisan tak menurun dari fungsi-fungsi terintegralkan yang konvergen
hampir di mana-mana ke $J\times g\phi$; menurut Teorema B.Levi,

\Centerline{$\int J\times g\phi\,d\mu$ ada $=\lim_{n\to\infty}
\int J\times g_n\phi\,d\mu=\lim_{n\to\infty}\int g_nd\nu
=\int g\,d\nu$.}

\medskip

{\bf (c)} Jika $g=g^+-g^-$, dengan $g^+$ dan $g^-$ fungsi-fungsi yang
$\nu$-terintegralkan, maka

\Centerline{$\int J\times g\phi\,d\mu
=\int J\times g^+\phi\,d\mu-\int J\times g^-\phi\,d\mu
=\int g^+d\nu-\int g^-d\nu
=\int g\,d\nu$.}

\medskip

{\bf (d)} Ini menangani kasus $H=Y$. Untuk kasus umum, kita mempunyai

$$\eqalignno{\int_Hg\,d\nu
&=\int(g\times\chi H)d\nu\cr
\displaycause{131Fa}
&=\int J\times(g\times\chi H)\phi\,d\mu
=\int J\times g\phi\times\chi(\phi^{-1}[H])d\mu
=\int_{\phi^{-1}[H]}J\times g\phi\,d\mu\cr}$$

\noindent menurut 214F.

\medskip

{\bf (e)} Untuk integral atas dan bawah, pertama-tama saya perhatikan
bahwa jika $F$ terabaikan terhadap $\nu$, maka
$\int J\times\chi(\phi^{-1}[F])d\mu=0$, sehingga
$J=0\,\,\mu$-hampir di mana-mana\ pada $\phi^{-1}[F]$. Akibatnya, jika
$f$ dan $g$ merupakan fungsi bernilai $[-\infty,\infty]$ pada
subhimpunan-subhimpunan $Y$ dan $f\leae g$, maka
$J\times f\phi\leae J\times g\phi$. Sekarang, jika
$\overline{\int}fd\nu=\infty$, tentu kita mempunyai
$\overline{\int}J\times f\phi\,d\mu\le\overline{\int}fd\nu$. Jika tidak,

$$\eqalign{\overline{\int}fd\nu
&=\inf\{\int g\,d\nu:g\text{ terintegralkan terhadap }\nu
  \text{ dan }f\leae g\}\cr
&=\inf\{\int J\times g\phi\,d\mu:
  g\text{ terintegralkan terhadap }\nu\text{ dan }f\leae g\}\cr
&\ge\inf\{\int h\,d\mu:
  h\text{ terintegralkan terhadap }\mu\text{ dan }J\times f\phi\leae h\}
=\overline{\intop}J\times f\phi\,d\mu.\cr}$$

\noindent Demikian pula, atau dengan menerapkan argumen ini pada $-f$,
kita memperoleh
$\underline{\int}J\times f\phi\,d\mu\ge\underline{\int}f\,d\nu$.
}%end of proof of 235A

\cmmnt{
\leader{235B}{Catatan (a)} Perhatikan konvensi khusus

\Centerline{$0\times\text{tidak terdefinisi}=0$}

\noindent yang saya terapkan dalam menafsirkan $J\times g\phi$. Ini
merupakan pokok teknis pertama dari sejumlah pokok yang akan kita hadapi
dalam bagian ini. Pokoknya ialah bahwa jika $g$ terdefinisi
$\nu$-hampir di mana-mana, maka untuk setiap perluasan $g$ menjadi fungsi
$g_1:Y\to\Bbb R$, berdasarkan konvensi ini kita akan mempunyai
$J\times g\phi=J\times g_1\phi$ kecuali pada
$\{x:J(x)>0,\,\phi(x)\in Y\setminus\dom g\}$, yang terabaikan; sehingga

\Centerline{$\int J\times g\phi\,d\mu
=\int J\times g_1\phi\,d\mu
=\int g_1d\nu
=\int g\,d\nu$}

\noindent jika $g$ dan $g_1$ terintegralkan. Dengan demikian, konvensi
tersebut sesuai di sini. Meskipun menambahkan satu frasa pada pernyataan
banyak hasil dalam bagian ini, konvensi tersebut memperlancar
penerapannya. (Tetapi saya harus menegaskan bahwa saya menggunakannya
hanya sebagai konvensi lokal, dan kaidah biasa
$0\times\text{tidak terdefinisi}=\text{tidak terdefinisi}$ akan tetap
berlaku di tempat lain dalam risalah ini kecuali jika secara eksplisit
dikesampingkan.)

\header{235Bb}{\bf (b)} Dalam merumuskan teorema ini, saya harus
berhati-hati membedakan hipotesis

\Centerline{$\int J(x)\chi(\phi^{-1}[F])(x)\mu(dx)$ ada $=\nu F$
setiap kali $\nu F<\infty$}

\noindent dari alternatif yang mungkin lebih elok

\Centerline{$\int_{\phi^{-1}[F]}J(x)\mu(dx)$ ada $=\nu F$ setiap kali
$\nu F<\infty$,}

\noindent yang tidak sepenuhnya memadai bagi teorema tersebut. (Lihat
235Q di bawah.) Ingat bahwa dengan $\int_Af$ saya maksud
$\int(f\restr A)d\mu_A$, dengan $\mu_A$ ukuran subruang pada $A$ (214D).
Mungkin saja $\int_A(f\restr A)d\mu_A$ terdefinisi bahkan ketika
$\int f\times\chi A\,d\mu$ tidak terdefinisi; misalnya, ambil $\mu$
sebagai ukuran Lebesgue pada $[0,1]$, $A$ sembarang subhimpunan tak
terukur dari $[0,1]$, dan $f$ fungsi konstan bernilai $1$; maka
$\int_Af=\mu^*A$, tetapi $f\times\chi A=\chi A$ tidak
$\mu$-terintegralkan. Namun, jika $\int f\times\chi A\,d\mu$
terdefinisi, maka $\int_Af$ juga terdefinisi dan keduanya sama; ini
merupakan akibat 214F. Walaupun 235P menunjukkan bahwa pada kebanyakan
kasus yang relevan bagi jilid ini perbedaan tersebut dapat dilewati,
penting untuk tidak mengandaikan bahwa $\phi^{-1}[F]$ terukur bagi setiap
$F\in\Tau$. Contoh sederhananya adalah sebagai berikut. Tetapkan
$X=Y=[0,1]$. Misalkan $\mu$ ukuran Lebesgue pada $[0,1]$ dan definisikan
$\nu$ dengan menetapkan

\Centerline{$\Tau=\{F:F\subseteq[0,1],\,F\cap[0,\bover12]$ terukur
Lebesgue$\}$,}

\Centerline{$\nu F=2\mu(F\cap[0,\bover12])$ untuk setiap $F\in\Tau$.}

\noindent Tetapkan $\phi(x)=x$ untuk setiap $x\in[0,1]$. Maka kita
mempunyai

\Centerline{$\nu F=\int_FJ\,d\mu
=\int J\times\chi(\phi^{-1}[F])d\mu$}

\noindent untuk setiap $F\in\Tau$, dengan $J(x)=2$ untuk
$x\in[0,\bover12]$ dan $J(x)=0$ untuk
$x\in\ocint{\bover12,1}$. Namun tentu terdapat subhimpunan $F$ dari
$[\bover12,1]$ yang tidak terukur Lebesgue (lihat 134D), dan $F$ semacam
itu mesti termasuk dalam $\Tau$, walaupun $\phi^{-1}[F]$ tidak termasuk
dalam domain $\Sigma$ dari $\mu$.

Pokoknya di sini ialah bahwa jika $\nu F_0=0$, kita mengharapkan
$J=0$ pada $\phi^{-1}[F_0]$, dan tidak penting apakah
$\phi^{-1}[F]$ terukur bagi $F\subseteq F_0$.
}%end of comment

\leader{235C}{}\cmmnt{ Teorema 235A menyangkut integrasi, sehingga
hipotesis $\int J\times\chi(\phi^{-1}[F])d\mu=\nu F$ hanya melihat
himpunan $F$ yang berukuran hingga. Jika kita hendak meninjau
keterukuran fungsi yang tidak terintegralkan, kita memerlukan hipotesis
yang sedikit lebih kuat. Saya mendekati versi ini secara lebih bertahap,
dengan sepasang lema.

\medskip

\noindent}{\bf Lema} Misalkan $\Sigma$ dan $\Tau$ aljabar-$\sigma$ dari
subhimpunan-subhimpunan $X$ dan $Y$ berturut-turut. Andaikan
$D\subseteq X$ dan $\phi:D\to Y$ suatu fungsi sedemikian sehingga
$\phi^{-1}[F]\in\Sigma_D$, aljabar-$\sigma$ subruang, untuk setiap
$F\in\Tau$. Maka $g\phi$ terukur terhadap $\Sigma$ bagi setiap fungsi
bernilai $[-\infty,\infty]$ yang terukur terhadap $\Tau$ dan dinotasikan
dengan $g$, serta didefinisikan pada suatu subhimpunan $Y$.

\proof{ Tetapkan $C=\dom g$ dan $B=\dom g\phi=\phi^{-1}[C]$. Jika
$a\in\Bbb R$, maka terdapat $F\in\Tau$ sedemikian sehingga
$\{y:g(y)\le a\}=F\cap C$. Sekarang terdapat $E\in\Sigma$ sedemikian
sehingga $\phi^{-1}[F]=E\cap D$. Jadi

\Centerline{$\{x:g\phi(x)\le a\}=B\cap E\in\Sigma_B$.}

\noindent Karena $a$ sembarang, $g\phi$ terukur terhadap $\Sigma$.
}%end of proof of 235C

\leader{235D}{}\cmmnt{ Sebagian hasil di bawah lebih mudah bila kita
dapat berpindah dengan bebas antara ruang ukur dan pelengkapannya (212C).
Lema berikut adalah yang kita perlukan.

\medskip

\noindent}{\bf Lema} Misalkan $(X,\Sigma,\mu)$ dan $(Y,\Tau,\nu)$
ruang-ruang ukur, dengan pelengkapan $(X,\hat\Sigma,\hat\mu)$ dan
$(Y,\hat\Tau,\hat\nu)$. Misalkan $\phi:D_{\phi}\to Y$ dan
$J:D_J\to\coint{0,\infty}$ fungsi-fungsi yang didefinisikan pada
subhimpunan koterabaikan dari $X$.

(a) Jika $\int J\times\chi(\phi^{-1}[F])d\mu=\nu F$ setiap kali
$F\in\Tau$ dan $\nu F<\infty$, maka
$\int J\times\chi(\phi^{-1}[F])d\hat\mu=\hat\nu F$ setiap kali
$F\in\hat\Tau$ dan $\hat\nu F<\infty$.

(b) Jika $\int J\times\chi(\phi^{-1}[F])d\mu=\nu F$ setiap kali
$F\in\Tau$, maka
$\int J\times\chi(\phi^{-1}[F])d\hat\mu=\hat\nu F$ setiap kali
$F\in\hat\Tau$.

\proof{ Keduanya bertumpu pada fakta bahwa salah satu hipotesis cukup
untuk memastikan $\int J\times\chi(\phi^{-1}[F])d\mu=0$ setiap kali
$\nu F=0$. Oleh karena itu, jika $F$ terabaikan terhadap $\nu$, sehingga
terdapat $F'\in\Tau$ dengan $F\subseteq F'$ dan $\nu F'=0$, kita akan
mempunyai

\Centerline{$\int J\times\chi(\phi^{-1}[F])d\mu
=\int J\times\chi(\phi^{-1}[F'])d\mu=0$.}

\noindent Namun sekarang, diberikan sembarang $F\in\hat\Tau$, terdapat
$F_0\in\Tau$ sedemikian sehingga $F_0\subseteq F$ dan
$\hat\nu(F\setminus F_0)=0$, sehingga

$$\eqalign{\int J\times\chi(\phi^{-1}[F])d\hat\mu
&=\int J\times\chi(\phi^{-1}[F])d\mu\cr
&=\int J\times\chi(\phi^{-1}[F_0])d\mu
   +\int J\times\chi(\phi^{-1}[F\setminus F_0])d\mu\cr
&=\nu F_0=\hat\nu F,\cr}$$

\noindent dengan ketentuan (untuk bagian (a)) bahwa
$\hat\nu F<\infty$.
}%end of proof of 235D

\cmmnt{\medskip

\noindent{\bf Catatan} Jadi, jika hipotesis salah satu hasil utama
bagian ini berlaku bagi sepasang ruang ukur yang tidak lengkap, kita
dapat mengharapkan suatu kesimpulan dengan menerapkan hasil tersebut
pada pelengkapan ruang-ruang yang diberikan.
}%end of comment

\leader{235E}{}\cmmnt{ Sekarang saya sampai pada versi alternatif 235A.

\medskip

\noindent}{\bf Teorema} Misalkan $(X,\Sigma,\mu)$ dan $(Y,\Tau,\nu)$
ruang-ruang ukur, serta $\phi:D_{\phi}\to Y$ dan
$J:D_J\to\coint{0,\infty}$ dua fungsi yang didefinisikan pada subhimpunan
koterabaikan dari $X$ sedemikian sehingga

\Centerline{$\int J\times\chi(\phi^{-1}[F])d\mu=\nu F$}

\noindent untuk setiap $F\in\Tau$, dengan memperbolehkan $\infty$ sebagai
nilai integral.

(a) $J\times g\phi$ terukur secara virtual terhadap $\mu$ bagi setiap
fungsi yang terukur secara virtual terhadap $\nu$ dan dinotasikan dengan
$g$, serta didefinisikan pada suatu subhimpunan $Y$.

(b) Misalkan $g$ suatu fungsi yang terukur secara virtual terhadap
$\nu$, bernilai $[-\infty,\infty]$, dan didefinisikan pada subhimpunan
koterabaikan dari $Y$. Maka
$\int J\times g\phi\,d\mu=\int g\,d\nu$ setiap kali salah satu integral
terdefinisi dalam $[-\infty,\infty]$, jika kita menafsirkan
$(J\times g\phi)(x)$ sebagai $0$ ketika $J(x)=0$ dan $g(\phi(x))$ tidak
terdefinisi.

\proof{ Misalkan $(X,\hat\Sigma,\hat\mu)$ dan $(Y,\hat\Tau,\hat\nu)$
pelengkapan dari $(X,\Sigma,\mu)$ dan $(Y,\Tau,\nu)$. Menurut 235D,

\Centerline{$\int J\times\chi(\phi^{-1}[F])d\hat\mu=\hat\nu F$}

\noindent untuk setiap $F\in\hat\Tau$. Dengan mengingat bahwa suatu
fungsi bernilai real terukur secara virtual terhadap $\mu$ jika dan hanya
jika terukur terhadap $\hat\Sigma$ (212Fa), dan bahwa
$\int fd\mu=\int fd\hat\mu$ jika salah satunya terdefinisi dalam
$[-\infty,\infty]$ (212Fb), kesimpulan yang kita cari adalah

\inset{(a)$'\,\,J\times g\phi$ terukur terhadap $\hat\Sigma$ bagi setiap
fungsi yang terukur terhadap $\hat\Tau$ dan dinotasikan dengan $g$, serta
didefinisikan pada suatu subhimpunan $Y$;}

\inset{(b)$'\,\,\int J\times g\phi\,d\hat\mu=\int g\,d\hat\nu$ setiap
kali $g$ suatu fungsi terukur terhadap $\hat\Tau$ yang didefinisikan
hampir di mana-mana dalam $Y$ dan salah satu integral terdefinisi dalam
$[-\infty,\infty]$.}

\medskip

{\bf (a)} Ketika saya menulis

\Centerline{$\int J\times\chi D_{\phi}d\mu
=\int J\times\chi(\phi^{-1}[Y])d\mu=\nu Y$,}

\noindent yang merupakan bagian dari hipotesis teorema ini, saya
bermaksud menyiratkan bahwa $J\times\chi D_{\phi}$ terukur secara virtual
terhadap $\mu$, yakni terukur terhadap $\hat\Sigma$. Karena $D_{\phi}$
koterabaikan, diperoleh bahwa $J$ terukur terhadap $\hat\Sigma$, dan
domainnya $D_J$, karena koterabaikan, juga termasuk dalam
$\hat\Sigma$. Tetapkan
$G=\{x:x\in D_J,\,J(x)>0\}\in\hat\Sigma$. Maka untuk setiap
$A\subseteq X$, $J\times\chi A$ terukur terhadap $\hat\Sigma$ jika dan
hanya jika $A\cap G\in\hat\Sigma$. Jadi, hipotesis tersebut tepat berarti
bahwa $G\cap\phi^{-1}[F]\in\hat\Sigma$ untuk setiap $F\in\hat\Tau$.

Sekarang misalkan $g$ fungsi bernilai $[-\infty,\infty]$, yang
didefinisikan pada suatu subhimpunan $C$ dari $Y$ dan terukur terhadap
$\hat\Tau$. Dengan menerapkan 235C pada $\phi\restr G$, kita melihat
bahwa $g\phi\restr G$ terukur terhadap $\hat\Sigma$, sehingga
$(J\times g\phi)\restr G$ terukur terhadap $\hat\Sigma$. Di pihak lain,
$J\times g\phi$ bernilai nol hampir di mana-mana dalam $X\setminus G$,
sehingga (karena $G\in\hat\Sigma$) $J\times g\phi$ terukur terhadap
$\hat\Sigma$, sebagaimana diperlukan.

\medskip

{\bf (b)(i)} Pertama-tama andaikan $g\ge 0$. Maka
$J\times g\phi\ge 0$, sehingga (a) menyatakan bahwa
$\int J\times g\phi$ terdefinisi dalam $[0,\infty]$.

\medskip

\qquad\grheada\ Jika $\int g\,d\hat\nu<\infty$, maka
$\int J\times g\phi\,d\hat\mu=\int g \,d\hat\nu$ menurut 235A.

\medskip

\qquad\grheadb\ Jika terdapat suatu $\epsilon>0$ sedemikian sehingga
$\hat\nu H=\infty$, dengan $H=\{y:g(y)\ge\epsilon\}$, maka

\Centerline{$\int J\times g\phi\,d\hat\mu
\ge\epsilon\int J\times\chi(\phi^{-1}[H])d\hat\mu
=\epsilon\hat\nu H=\infty$,}

\noindent sehingga

\Centerline{$\int J\times g\phi\,d\hat\mu=\infty=\int
g\,d\hat\nu$.}

\medskip

\qquad\grheadc\ Jika tidak,

$$\eqalign{\int J\times g\phi\,d\hat\mu
&\ge\sup\{\int J\times h\phi\,d\hat\mu:
   h\text{ adalah }\hat\nu\text{-terintegralkan},
   \,0\le h\le g\}\cr
&=\sup\{\int h\,d\hat\nu:h\text{ adalah }\hat\nu\text{-terintegralkan},
   \,0\le h\le g\}
=\int g\,d\hat\nu
=\infty,\cr}$$

\noindent sehingga sekali lagi
$\int J\times g\phi\,d\hat\mu=\int g\,d\hat\nu$.

\medskip

\quad{\bf (ii)} Untuk $g$ bernilai real umum, terapkan (i) pada $g^+$
dan $g^-$, dengan $g^+=\bover12(|g|+g)$,
$g^-=\bover12(|g|-g)$; pokoknya ialah bahwa
$(J\times g\phi)^+=J\times g^+\phi$ dan
$(J\times g\phi)^-=J\times g^-\phi$, sehingga

\Centerline{$\int J\times g\phi
=\int J\times g^+\phi-\int J\times g^-\phi
=\int g^+-\int g^- %
=\int g$}

\noindent jika salah satu ruas terdefinisi dalam $[-\infty,\infty]$.
}%end of proof of 235E

\cmmnt{
\leader{235F}{Catatan (a)} Tentu terdapat dua kasus khusus teorema ini
yang bersama-sama memuat seluruh isinya: kasus $J=1$ hampir di mana-mana\
dan kasus ketika $X=Y$ serta $\phi$ merupakan fungsi identitas. Jika
$J=\chi X$, kita sangat dekat dengan 235G di bawah, dan jika $\phi$
merupakan fungsi identitas, kita dekat dengan ukuran-ukuran integral tak
tentu dalam \S234.

\header{235Fb}{\bf (b)} Seperti dalam 235A, kita dapat memperkuat
kesimpulan (b) dalam 235E menjadi

\Centerline{$\int_{\phi^{-1}[F]}J\times
g\phi\,d\mu=\int_Fg\,d\nu$}

\noindent setiap kali $F\in\Tau$ dan $\int_Fg\,d\nu$ terdefinisi dalam
$[-\infty,\infty]$.
}%end of comment

\leader{235G}{Teorema} Misalkan $(X,\Sigma,\mu)$ dan $(Y,\Tau,\nu)$
ruang-ruang ukur dan $\phi:X\to Y$ suatu fungsi \imp. Maka

(a) jika $g$ suatu fungsi yang terukur secara virtual terhadap $\nu$,
bernilai $[-\infty,\infty]$, dan didefinisikan pada suatu subhimpunan
$Y$, maka $g\phi$ terukur secara virtual terhadap $\mu$;

(b) jika $g$ suatu fungsi yang terukur secara virtual terhadap $\nu$,
bernilai $[-\infty,\infty]$, dan didefinisikan pada suatu subhimpunan
koterabaikan dari $Y$, maka $\int g\phi\,d\mu=\int g\,d\nu$ jika salah
satu integral terdefinisi dalam $[-\infty,\infty]$;

(c) jika $g$ suatu fungsi yang terukur secara virtual terhadap $\nu$,
bernilai $[-\infty,\infty]$, dan didefinisikan pada suatu subhimpunan
koterabaikan dari $Y$, serta $F\in\Tau$, maka
$\int_{\phi^{-1}[F]}g\phi\,d\mu=\int_Fg\,d\nu$ jika salah satu integral
terdefinisi dalam $[-\infty,\infty]$.

\proof{{\bf (a)} Ini langsung mengikuti 234Ba dan 235C; dengan mengambil
$\hat\Sigma$ dan $\hat\Tau$ sebagai domain pelengkapan $\mu$ dan $\nu$
berturut-turut, $\phi^{-1}[F]\in\hat\Sigma$ untuk setiap
$F\in\hat\Tau$, sehingga jika $g$ terukur terhadap $\hat\Tau$, maka
$g\phi$ terukur terhadap $\hat\Sigma$.

\medskip

{\bf (b)} Terapkan 235E dengan $J=\chi X$; kita mempunyai

\Centerline{$\int J\times\chi(\phi^{-1}[F])d\mu
=\mu\phi^{-1}[F]=\nu F$}

\noindent untuk setiap $F\in\Tau$, sehingga

\Centerline{$\int g\phi=\int J\times g\phi=\int g$}

\noindent jika salah satu integral terdefinisi dalam $[-\infty,\infty]$.

\medskip

{\bf (c)} Terapkan (b) pada $g\times\chi F$.
}%end of proof of 235G

\cmmnt{
\leader{235H}{Malapetaka ukuran citra}
Penerapan 235A akan berlangsung jauh lebih lancar jika kita dapat
mengatakan

\inset{`$\int g\,d\nu$ ada dan sama dengan
$\int J\times g\phi\,d\mu$
untuk setiap $g:Y\to\Bbb R$ sedemikian sehingga $J\times g\phi$
$\mu$-terintegralkan'.}

\noindent Sayangnya, tidak ada harapan bagi suatu hasil yang dapat
diterapkan secara universal ke arah ini. Misalkan, misalnya, $\nu$
ukuran Lebesgue pada $Y=[0,1]$, $X\subseteq[0,1]$ suatu himpunan yang
tidak terukur Lebesgue dan berukuran luar $1$ (134D), $\mu$ ukuran
subruang $\nu_X$ pada $X$, dan $\phi(x)=x$ untuk $x\in X$. Maka

\Centerline{$\mu\phi^{-1}F=\nu^*(X\cap F)=\nu F$}

\noindent untuk setiap himpunan terukur Lebesgue $F\subseteq Y$, sehingga
kita dapat mengambil $J=\chi X$ dan hipotesis 235A serta 235E akan
terpenuhi. Namun, jika kita menuliskan $g=\chi X:[0,1]\to\{0,1\}$, maka
$\int g\phi\,d\mu$ terdefinisi meskipun $\int g\,d\nu$ tidak terdefinisi.

Pokoknya di sini ialah bahwa terdapat himpunan $A\subseteq Y$ sedemikian
sehingga (dalam bahasa 235A/235E) $\phi^{-1}[A]\in\Sigma$, tetapi
$A\notin\hat\Tau$. Inilah {\bf malapetaka ukuran citra}.
Pencarian konteks yang memungkinkan kita memastikan bahwa malapetaka ini
tidak terjadi akan menjadi salah satu tema penggerak Jilid 4. Untuk saat
ini, saya akan menawarkan beberapa catatan umum (235I-235J), lalu
menjelaskan salah satu kasus penting yang bebas dari masalah ini (235K).
}%end of comment


\leader{235I}{Lema}
Misalkan $\Sigma$ dan $\Tau$ aljabar-$\sigma$ dari
subhimpunan-subhimpunan $X$ dan $Y$ berturut-turut, serta $\phi$ suatu
fungsi dari suatu subhimpunan $D$ dari $X$ ke $Y$. Andaikan
$G\subseteq X$ dan

\Centerline{$\Tau=\{F:F\subseteq Y,\,G\cap\phi^{-1}[F]\in\Sigma\}$.}

\noindent Maka suatu fungsi bernilai real $g$, yang didefinisikan pada
suatu anggota $\Tau$, terukur terhadap $\Tau$ jika dan hanya jika
$\chi G\times g\phi$ terukur terhadap $\Sigma$.

\proof{ Karena tentu $Y\in\Tau$, hipotesis tersebut menyiratkan bahwa
$G\cap D=G\cap\phi^{-1}[Y]$ termasuk dalam $\Sigma$.

Misalkan $g:C\to\Bbb R$ suatu fungsi, dengan $C\in\Tau$. Tetapkan
$B=\dom(g\phi)=\phi^{-1}[C]$, dan untuk $a\in\Bbb R$ tetapkan
$F_a=\{y:g(y)\ge a\}$,

\Centerline{$E_a=G\cap\phi^{-1}[F_a]
=\{x:x\in G\cap B,\,g\phi(x)\ge a\}$.}

\noindent Perhatikan bahwa $G\cap B\in\Sigma$ karena $C\in\Tau$.

\medskip

\quad{\bf (i)} Jika $g$ terukur terhadap $\Tau$, maka $F_a\in\Tau$ dan
$E_a\in\Sigma$ untuk setiap $a$. Sekarang

\Centerline{$G\cap\{x:x\in B,\,g\phi(x)\ge a\}=G\cap\phi^{-1}[F_a]
=E_a$,}

\noindent sehingga $\{x:x\in B,\,(\chi G\times g\phi)(x)\ge a\}$ sama
dengan $E_a$ atau $E_a\cup(B\setminus G)$, dan dalam kedua kasus terukur
relatif terhadap $\Sigma$ dalam $B$. Karena $a$ sembarang,
$\chi G\times g\phi$ terukur terhadap $\Sigma$.

\medskip

\quad{\bf (ii)} Jika $\chi G\times g\phi$ terukur terhadap $\Sigma$,
maka untuk setiap $a\in\Bbb R$,

\Centerline{$E_a=\{x:x\in G\cap B,\,(\chi G\times g\phi)(x)\ge
a\}\in\Sigma$}

\noindent karena $G\cap B\in\Sigma$ dan $\chi G\times g\phi$ terukur
terhadap $\Sigma$. Jadi $F_a\in\Tau$. Karena $a$ sembarang, $g$ terukur
terhadap $\Tau$.
}%end of proof of 235I

\leader{235J}{Teorema}
Misalkan $(X,\Sigma,\mu)$ dan $(Y,\Tau,\nu)$ ruang-ruang ukur lengkap.
Misalkan $\phi:D_{\phi}\to Y$ dan $J:D_J\to\coint{0,\infty}$
fungsi-fungsi yang didefinisikan pada subhimpunan koterabaikan dari $X$,
lalu tetapkan $G=\{x:x\in D_J,\,J(x)>0\}$. Andaikan

\Centerline{$\Tau
=\{F:F\subseteq Y,\,G\cap\phi^{-1}[F]\in\Sigma\}$,}

\Centerline{$\nu F=\int J\times\chi(\phi^{-1}[F])d\mu$ untuk setiap
$F\in\Tau$.}

\noindent Maka, untuk setiap fungsi bernilai real $g$ yang didefinisikan
pada suatu subhimpunan $Y$,
$\int J\times g\phi\,d\mu=\int g\,d\nu$ setiap kali salah satu integral
terdefinisi dalam $[-\infty,\infty]$, dengan ketentuan bahwa kita
menafsirkan $(J\times g\phi)(x)$ sebagai $0$ ketika $J(x)=0$ dan
$g(\phi(x))$ tidak terdefinisi.

\proof{ Jika $g$ terukur terhadap $\Tau$ dan didefinisikan hampir di
mana-mana, ini merupakan akibat 235E. Jadi, saya harus menunjukkan bahwa
jika $J\times g\phi$ terukur dan didefinisikan hampir di mana-mana, maka
$g$ juga demikian. Tetapkan $W=Y\setminus\dom g$. Maka
$J\times g\phi$ tidak terdefinisi pada $G\cap\phi^{-1}[W]$, sebab
$g\phi$ tidak terdefinisi di sana dan kita tidak dapat memanfaatkan
klausul pelepasan yang tersedia ketika $J=0$; jadi
$G\cap\phi^{-1}[W]$ mesti terabaikan, sehingga terukur, dan $W\in\Tau$.
Selanjutnya,

\Centerline{$\nu W=\int J\times\chi(\phi^{-1}[W])=0$}

\noindent karena $J\times\chi(\phi^{-1}[W])$ hanya dapat bernilai tak nol
pada himpunan terabaikan $G\cap\phi^{-1}[W]$. Jadi $g$ terdefinisi hampir
di mana-mana.

Perhatikan bahwa hipotesis tentu menyiratkan bahwa $J\times\chi
D_{\phi}=J\times\chi(\phi^{-1}[Y])$ terukur, sehingga $J$ terukur
(karena $D_{\phi}$ koterabaikan) dan $G\in\Sigma$. Dengan menuliskan
$K(x)=1/J(x)$ untuk $x\in G$ dan $0$ untuk $x\in X\setminus G$, fungsi
$K:X\to\Bbb R$ terukur, dan

\Centerline{$\chi G\times g\phi=K\times J\times g\phi$}

\noindent terukur. Jadi, 235I menyatakan bahwa $g$ mesti terukur, dan
selesailah bukti.
}%end of proof of 235J

\cmmnt{\medskip

\noindent{\bf Catatan} Ketika $J=\chi X$, hipotesis teorema ini menjadi

\Centerline{$\Tau=\{F:F\subseteq Y,\,\phi^{-1}[F]\in\Sigma\}$,
\quad $\nu F=\mu\phi^{-1}[F]$ untuk setiap $F\in\Tau$;}

\noindent yakni, $\nu$ merupakan ukuran citra $\mu\phi^{-1}$ sebagaimana
didefinisikan dalam 234D.
}%end of comment

\leader{235K}{Akibat}
Misalkan $(X,\Sigma,\mu)$ suatu ruang ukur
lengkap, dan $J$ suatu fungsi terukur nonnegatif yang didefinisikan pada
suatu himpunan bagian koterabaikan dari $X$. Misalkan $\nu$ merupakan
ukuran integral tak tentu terkait, dan $\Tau$ domainnya. Maka untuk setiap
fungsi bernilai real $g$ yang didefinisikan pada suatu himpunan bagian
dari $X$, $g$ terukur terhadap $\Tau$ jika dan hanya jika $J\times g$
terukur terhadap $\Sigma$, dan
$\int g\,d\nu=\int J\times g\,d\mu$ jika salah satu integral terdefinisi
dalam $[-\infty,\infty]$, asalkan $(J\times g)(x)$ kita tafsirkan sebagai
$0$ ketika $J(x)=0$ dan $g(x)$ tidak terdefinisi.

\proof{ Gabungkan 235J, dengan mengambil $Y=X$ dan $\phi$ fungsi identitas,
dengan 234Ld.
}%end of proof of 235K

\cmmnt{
\leader{235L}{Menerapkan Teorema Radon-Nikod\'ym}
Agar dapat
menggunakan 235A--235J secara efektif, kita perlu dapat menemukan fungsi
$J$ yang sesuai. Hal ini dapat sulit dilakukan -- beberapa contoh yang
sangat khusus akan mengisi sebagian besar Bab 26 di bawah. Namun terdapat
banyak keadaan yang menjamin keberadaan $J$ demikian, sekalipun kita tidak
mengetahui bentuknya. Syarat minimalnya ialah bahwa jika $\nu F<\infty$
dan $\mu^*\phi^{-1}[F]=0$ maka $\nu F=0$, karena
$\int J\times\chi(\phi^{-1}[F])d\mu$ akan bernilai nol untuk setiap $J$.
Suatu syarat cukup, dalam kasus khusus ukuran integral tak tentu, diberikan
dalam 234O. Syarat cukup lainnya adalah yang berikut.
}%end of comment

\leader{235M}{Teorema}
Misalkan $(X,\Sigma,\mu)$ suatu
ruang ukur $\sigma$-hingga, $(Y,\Tau,\nu)$ suatu ruang ukur semihingga,
dan $\phi:D\to Y$ suatu fungsi sedemikian sehingga

(i) $D$ merupakan himpunan bagian koterabaikan dari $X$,

(ii) $\phi^{-1}[F]\in\Sigma$ untuk setiap $F\in\Tau$;

(iii) $\mu\phi^{-1}[F]>0$ setiap kali $F\in\Tau$ dan $\nu F>0$.

\noindent Maka terdapat fungsi yang terukur terhadap $\Sigma$, yaitu
$J:X\to\coint{0,\infty}$, sedemikian sehingga
$\int J\times\chi\phi^{-1}[F]\,d\mu=\nu F$ untuk
setiap $F\in\Tau$.

\proof{{\bf (a)} Sebagai permulaan (hingga akhir bagian (c) di bawah),
anggaplah bahwa $D=X$ dan bahwa $\nu$ berhingga total.

Tetapkan $\tildeTau=\{\phi^{-1}[F]:F\in\Tau\}\subseteq\Sigma$. Maka
$\tildeTau$ merupakan aljabar-$\sigma$ dari himpunan-himpunan bagian $X$.
\Prf\ (i)

\Centerline{$\emptyset=\phi^{-1}[\emptyset]\in\tildeTau$.}

\noindent (ii) Jika $E\in\tildeTau$, ambil $F\in\Tau$ sedemikian sehingga
$E=\phi^{-1}[F]$, sehingga

\Centerline{$X\setminus E=\phi^{-1}[Y\setminus F]\in\tildeTau$.}

\noindent (iii) Jika $\sequencen{E_n}$ sembarang barisan dalam
$\tildeTau$, maka untuk setiap $n\in\Bbb N$ pilih $F_n\in\Tau$ sedemikian
sehingga $E_n=\phi^{-1}[F_n]$; maka

\Centerline{$\bigcup_{n\in\Bbb N}E_n=\phi^{-1}[\bigcup_{n\in\Bbb
N}F_n]\in\tildeTau$. \Qed}

\noindent Selanjutnya, kita mempunyai ukuran berhingga total
$\tilde\nu:\tildeTau\to[0,\nu Y]$ yang diberikan dengan menetapkan

\Centerline{$\tilde\nu(\phi^{-1}[F])=\nu F$ untuk setiap $F\in\Tau$.}

\noindent\Prf\ (i) Jika $F$, $F'\in\Tau$ dan
$\phi^{-1}[F]=\phi^{-1}[F']$, maka
$\phi^{-1}[F\symmdiff F']=\emptyset$, sehingga
$\mu(\phi^{-1}[F\symmdiff F'])=0$ dan $\nu(F\symmdiff F')=0$;
akibatnya $\nu F=\nu F'$. Ini menunjukkan bahwa $\tilde\nu$
terdefinisi dengan baik. (ii) Sekarang

\Centerline{$\tilde\nu\emptyset
=\tilde\nu(\phi^{-1}[\emptyset])
=\nu\emptyset=0$.}

\noindent (iii) Jika $\sequencen{E_n}$ suatu barisan saling lepas dalam
$\tildeTau$, ambil barisan $\sequencen{F_n}$ dalam $\Tau$ sedemikian
sehingga $E_n=\phi^{-1}[F_n]$ untuk setiap $n$; tetapkan
$F'_n=F_n\setminus\bigcup_{m<n}F_m$ untuk setiap $n$; maka
$E_n=\phi^{-1}[F'_n]$ untuk setiap $n$, sehingga

\Centerline{$\tilde\nu(\bigcup_{n\in\Bbb N}E_n)
=\tilde\nu(\phi^{-1}[\bigcup_{n\in\Bbb N}F'_n])
=\nu(\bigcup_{n\in\Bbb N}F'_n)
=\sum_{n=0}^{\infty}\nu F'_n
=\sum_{n=0}^{\infty}\tilde\nu E_n$. \Qed}

Terakhir, perhatikan bahwa jika $\tilde\nu E>0$ maka $\mu E>0$, karena
$E=\phi^{-1}[F]$ dengan $\nu F>0$.

\medskip

{\bf (b)} Menurut 215B(ix), terdapat fungsi yang terukur terhadap $\Sigma$,
yaitu $h:X\to\ooint{0,\infty}$, sedemikian sehingga $\int h\,d\mu$
berhingga.
Definisikan $\tilde\mu:\tildeTau\to\coint{0,\infty}$ dengan menetapkan
$\tilde\mu E=\int_Eh\,d\mu$ untuk setiap $E\in\tildeTau$; maka
$\tilde\mu$ merupakan ukuran berhingga total. Jika $E\in\tildeTau$ dan
$\tilde\mu E=0$, maka (karena $h$ positif ketat) $\mu E=0$ dan
$\tilde\nu E=0$. Dengan demikian, kita dapat menerapkan Teorema
Radon-Nikod\'ym pada $\tilde\mu$ dan $\tilde\nu$ untuk melihat bahwa
terdapat fungsi yang terukur terhadap $\tildeTau$, yaitu
$g:X\to\Bbb R$, sedemikian sehingga
$\int_Eg\,d\tilde\mu=\tilde\nu E$ untuk setiap
$E\in\tildeTau$. Karena $\tilde\nu$ nonnegatif, kita boleh menganggap
bahwa $g\ge 0$.

\medskip

{\bf (c)} Dengan menerapkan 235A pada $\mu$, $\tilde\mu$, $h$, dan
fungsi identitas dari $X$ ke dirinya sendiri, kita memperoleh

\Centerline{$\int_Eg\times h\,d\mu=\int_Eg\,d\tilde\mu
=\tilde\nu E$}

\noindent untuk setiap $E\in\tildeTau$, yakni

\Centerline{$\int J\times\chi(\phi^{-1}[F])d\mu = \nu F$}

\noindent untuk setiap $F\in\Tau$, dengan menuliskan $J=g\times h$.

\medskip

{\bf (d)} Ini menyelesaikan bukti ketika $\nu$ berhingga total dan
$D=X$. Untuk kasus umum, jika $Y=\emptyset$ maka $\mu X=0$ dan hasilnya
langsung. Jika tidak, misalkan $\hat\phi$ sembarang perluasan $\phi$
menjadi fungsi dari $X$ ke $Y$ yang konstan pada $X\setminus D$; maka
$\hat\phi^{-1}[F]\in\Sigma$ untuk setiap $F\in\Tau$, karena
$D=\phi^{-1}[Y]\in\Sigma$ dan $\hat\phi^{-1}[F]$ selalu salah satu dari
$\phi^{-1}[F]$ atau $(X\setminus D)\cup\phi^{-1}[F]$. Sekarang $\nu$
harus $\sigma$-hingga.
\Prf\ Gunakan kriteria 215B(ii). Jika $\Cal F$ suatu keluarga saling
lepas dalam $\{F:F\in\Tau,\,0<\nu F<\infty\}$, maka
$\Cal E=\{\hat\phi^{-1}[F]:F\in\Cal F\}$ merupakan keluarga saling lepas
dalam $\{E:\mu E>0\}$, sehingga $\Cal E$ dan $\Cal F$ terhitung. \Qed

Misalkan $\sequencen{Y_n}$ suatu partisi $Y$ menjadi himpunan-himpunan
dengan ukuran $\nu$ berhingga, dan untuk setiap $n\in\Bbb N$ tetapkan
$\nu_nF=\nu(F\cap Y_n)$ untuk setiap $F\in\Tau$. Maka $\nu_n$ merupakan
ukuran berhingga total pada $Y$, dan jika $\nu_nF>0$ maka $\nu F>0$
sehingga

\Centerline{$\mu\hat\phi^{-1}[F]=\mu\phi^{-1}[F]>0$.}

\noindent Dengan demikian $\mu$, $\hat\phi$, dan $\nu_n$ memenuhi
asumsi teorema bersama asumsi-asumsi bagian (a) di atas, dan terdapat
fungsi yang terukur terhadap $\Sigma$, yaitu
$J_n:X\to\coint{0,\infty}$, sedemikian sehingga

\Centerline{$\nu_nF=\int J_n\times\chi(\phi^{-1}[F])d\mu$}

\noindent untuk setiap $F\in\Tau$. Sekarang tetapkan
$J=\sum_{n=0}^{\infty}J_n\times\chi(\phi^{-1}[Y_n])$, sehingga
$J:X\to\coint{0,\infty}$ terukur terhadap $\Sigma$. Jika $F\in\Tau$,
maka

$$\eqalign{\int J\times\chi(\phi^{-1}[F])d\mu
&=\sum_{n=0}^{\infty}
  \int J_n\times\chi(\phi^{-1}[Y_n])\times\chi(\phi^{-1}[F])d\mu\cr
&=\sum_{n=0}^{\infty}\int J_n\times\chi(\phi^{-1}[F\cap Y_n])d\mu
=\sum_{n=0}^{\infty}\nu(F\cap Y_n)
=\nu F,\cr}$$

\noindent sebagaimana diperlukan.
}%end of proof of 235M

\leader{235N}{Catatan}
Teorema 235M dapat gagal jika $\mu$ hanya dapat
dilokalkan secara ketat alih-alih $\sigma$-hingga. \prooflet{\Prf\ Misalkan
$X=Y$ suatu himpunan tak terhitung, $\Sigma=\Cal PX$, $\mu$
ukuran pencacahan pada $X$ (112Bd), $\Tau$ aljabar-$\sigma$
terhitung-koterhitung dari $Y$, $\nu$ ukuran terhitung-koterhitung pada
$Y$ (211R), dan $\phi:X\to Y$ peta identitas. Maka
$\phi^{-1}[F]\in\Sigma$ dan $\mu\phi^{-1}[F]>0$ setiap kali $\nu F>0$.
Namun, jika $J$ sembarang fungsi terintegralkan terhadap $\mu$ pada $X$,
maka $F=\{x:J(x)\ne 0\}$ terhitung dan

\Centerline{$\nu(Y\setminus F)=1\ne 0
=\int_{\phi^{-1}[Y\setminus F]}J\,d\mu$. \Qed}
}

\leader{*235O}{}\cmmnt{ Terdapat beberapa penyederhanaan
dalam kasus ruang $\sigma$-hingga; khususnya, 235A dan 235E menjadi
menyatu. Saya akan memberikan suatu adaptasi hipotesis 235A yang dapat
digunakan dalam kasus $\sigma$-hingga. Pertama-tama sebuah lema.

\medskip

\noindent}{\bf Lema} Misalkan $(X,\Sigma,\mu)$ suatu ruang ukur dan $f$
suatu fungsi terintegralkan nonnegatif pada $X$. Jika $A\subseteq X$
sedemikian sehingga $\int_Af+\int_{X\setminus A}f=\int f$, maka
$f\times\chi A$ terintegralkan.

\proof{ Menurut 214Eb, terdapat fungsi-fungsi terintegralkan terhadap
$\mu$, $f_1$ dan $f_2$, sedemikian sehingga $f_1$ memperluas
$f\restr A$, $f_2$ memperluas $f\restr X\setminus A$, dan

\Centerline{$\int_Ef_1=\int_{E\cap A}f$,
\quad $\int_Ef_2=\int_{E\setminus A}f$}

\noindent untuk setiap $E\in\Sigma$. Karena $f$ nonnegatif,
$\int_Ef_1$ dan $\int_Ef_2$ nonnegatif untuk setiap $E\in\Sigma$, dan
$f_1$, $f_2$ nonnegatif hampir di mana-mana. Dengan demikian kita
mempunyai $f\times\chi A\leae f_1$ dan
$f\times\chi(X\setminus A)\leae f_2$, sehingga
$f\leae f_1+f_2$. Tetapi juga

\Centerline{$\int f_1+f_2
=\int_Xf_1+\int_Xf_2=\int_Af+\int_{X\setminus A}f=\int f$,}

\noindent sehingga $f\eae f_1+f_2$. Dengan demikian

\Centerline{$f_1\eae f-f_2\leae f-f\times\chi(X\setminus A)
=f\times\chi A\leae f_1$}

\noindent dan $f\times\chi A\eae f_1$ terintegralkan.
}%end of proof of 235O

\leader{*235P}{Proposisi}
Misalkan $(X,\Sigma,\mu)$ suatu ruang ukur
lengkap dan $(Y,\Tau,\nu)$ suatu ruang ukur $\sigma$-hingga lengkap.
Andaikan bahwa $\phi:D_{\phi}\to Y$,
$J:D_J\to\coint{0,\infty}$ merupakan fungsi-fungsi yang didefinisikan
pada himpunan-himpunan bagian koterabaikan $D_{\phi}$, $D_J$ dari $X$
sedemikian sehingga $\int_{\phi^{-1}[F]}J\,d\mu$ ada dan sama dengan
$\nu F$ setiap kali $F\in\Tau$ dan $\nu F<\infty$.

(a) $J\times g\phi$ terukur terhadap $\Sigma$ untuk setiap fungsi yang
terukur terhadap $\Tau$, bernilai real, dan dinotasikan dengan $g$, serta
didefinisikan pada suatu himpunan bagian dari $Y$.

(b) Jika $g$ suatu fungsi bernilai real terukur terhadap $\Tau$ yang
didefinisikan hampir di mana-mana dalam $Y$, maka
$\int J\times g\phi\,d\mu=\int g\,d\nu$ setiap kali salah satu integral
terdefinisi dalam $[-\infty,\infty]$, dengan menafsirkan
$(J\times g\phi)(x)$ sebagai $0$ ketika $J(x)=0$ dan $g(\phi(x))$ tidak
terdefinisi.

\proof{ Intinya adalah bahwa hipotesis-hipotesis 235E
terpenuhi. Untuk melihat ini, tuliskan
$\Sigma_C=\{E\cap C:E\in\Sigma\}$ dan
$\mu_C=\mu^*\restr \Sigma_C$ bagi ukuran subruang pada $C$, untuk setiap
$C\subseteq X$. Misalkan
$\sequencen{Y_n}$ suatu barisan tak menurun dari himpunan-himpunan dengan
gabungan $Y$ dan dengan $\nu Y_n<\infty$ untuk setiap $n\in\Bbb N$,
dimulai dari $Y_0=\emptyset$.

\medskip

{\bf (i)} Ambil sembarang $F\in\Tau$ dengan $\nu F<\infty$, dan
tetapkan $F_n=F\cup Y_n$ untuk setiap $n\in\Bbb N$; tuliskan
$C_n=\phi^{-1}[F_n]$.

Untuk sementara, tetapkan $n$. Hipotesis kita menyiratkan bahwa

\Centerline{$\int_{C_0} J\,d\mu+\int_{C_n\setminus C_0} J\,d\mu
=\nu F+\nu(F_n\setminus F)=\nu F_n=\int_{C_n}J\,d\mu$.}

\noindent Jika ukuran subruang pada $C_0$ dan $C_n\setminus C_0$ kita
pandang sebagai turunan dari ukuran $\mu_{C_n}$ pada $C_n$ (214Ce),
maka 235O memberi tahu kita bahwa $J\times\chi C_0$ terintegralkan
terhadap $\mu_{C_n}$, dan terdapat fungsi yang terintegralkan terhadap
$\mu$, yaitu $h_n$, sedemikian sehingga $h_n$ memperluas
$(J\times\chi C_0)\restr C_n$.

Misalkan $E$ suatu himpunan koterabaikan terhadap $\mu$, yang termuat
dalam domain $D_{\phi}$ dari $\phi$, sedemikian sehingga
$h_n\restr E$ terukur terhadap $\Sigma$ untuk setiap $n$. Karena
$\sequencen{C_n}$ merupakan barisan tak menurun dengan gabungan
$\phi^{-1}[\bigcup_{n\in\Bbb N}F_n]
\ifdim\pagewidth=390pt\penalty-100\fi
=D_{\phi}$,

\Centerline{$(J\times\chi C_0)(x)=\lim_{n\to\infty}h_n(x)$}

\noindent untuk setiap $x\in E$, dan $(J\times\chi C_0)\restr E$
terukur. Pada saat yang sama, kita mengetahui bahwa terdapat fungsi yang
terintegralkan terhadap $\mu$, yaitu $h$, yang memperluas $J\restr C_0$, dan
$0\leae J\times\chi C_0\leae|h|$. Dengan demikian
$J\times \chi C_0$ terintegralkan, dan (dengan menggunakan 214F)

\Centerline{$\int J\times\chi\phi^{-1}[F]\,d\mu
=\int J\times\chi C_0\,d\mu
=\int_{C_0}J\restr C_0\,d\mu_{C_0}
=\nu F$.}

\medskip

{\bf (ii)} Ini menangani $F$ berukuran berhingga. Untuk $F\in\Tau$
umum,

\Centerline{$\int J\times\chi(\phi^{-1}[F])\,d\mu
=\lim_{n\to\infty}\int J\times\chi(\phi^{-1}[F\cap Y_n])\,d\mu
=\lim_{n\to\infty}\nu(F\cap Y_n)
=\nu F$.}

\noindent Jadi hipotesis-hipotesis 235E terpenuhi, dan hasilnya segera
mengikuti.
}%end of proof of 235P

\leader{*235Q}{}\cmmnt{ Dalam 235Bb saya menyatakan bahwa suatu
kesulitan dapat muncul dalam 235A, untuk ruang ukur umum, jika kita
menuliskan $\int_{\phi^{-1}[F]}J\,d\mu$ dalam hipotesis sebagai
pengganti $\int J\times\chi(\phi^{-1}[F])d\mu$. Berikut sebuah contoh.
\medskip

\noindent}{\bf Contoh} Tetapkan $X=Y=[0,2]$. Tuliskan $\Sigma_L$ untuk
aljabar himpunan bagian terukur Lebesgue dari $\Bbb R$, dan $\mu_L$
untuk ukuran Lebesgue; tuliskan $\mu_c$ untuk ukuran pencacahan pada
$\Bbb R$. Tetapkan

\Centerline{$\Sigma=\Tau
=\{E:E\subseteq[0,2],\,E\cap\coint{0,1}\in\Sigma_L\}$\dvro{.}{;}}

\noindent\cmmnt{tentu saja ini merupakan aljabar-$\sigma$ dari himpunan-
himpunan bagian $[0,2]$. }Untuk $E\in\Sigma=\Tau$, tetapkan

\Centerline{$\mu E=\nu
E=\mu_L(E\cap\coint{0,1})+\mu_c(E\cap[1,2])$\dvro{.}{;}}

\noindent\cmmnt{maka $\mu$ merupakan ukuran lengkap -- pada dasarnya,
ia adalah jumlah langsung ukuran Lebesgue pada $\coint{0,1}$ dan ukuran
pencacahan pada $[1,2]$ (lihat 214L). Mudah dilihat bahwa

\Centerline{$\mu^*B=\mu_L^*(B\cap\coint{0,1})+\mu_c(B\cap[1,2])$}

\noindent untuk setiap $B\subseteq[0,2]$. }Misalkan
$A\subseteq\coint{0,1}$ suatu himpunan tak terukur Lebesgue sedemikian
sehingga $\mu_L^*(E\setminus A)=\mu_LE$ untuk setiap himpunan terukur
Lebesgue
$E\subseteq\coint{0,1}$\cmmnt{ (lihat 134D)}. Definisikan
$\phi:[0,2]\to[0,2]$ dengan
menetapkan $\phi(x)=x+1$ jika $x\in A$, dan $\phi(x)=x$ jika
$x\in[0,2]\setminus A$.

Jika $F\in\Sigma$, maka $\mu^*(\phi^{-1}[F])=\mu F$.
\prooflet{\Prf\
{\bf (i)} Jika
$F\cap[1,2]$ berhingga, maka
$\mu F=\mu_L(F\cap[0,1])+\#(F\cap[1,2])$.
Sekarang

\Centerline{$\phi^{-1}[F]=(F\cap\coint{0,1}\setminus
A)\cup(F\cap[1,2])\cup\{x:x\in A,\,x+1\in F\}$;}

\noindent karena himpunan terakhir berhingga, dan karena itu terabaikan
terhadap $\mu$,

\Centerline{$\mu^*(\phi^{-1}[F])
=\mu_L^*(F\cap\coint{0,1}\setminus A)+\#(F\cap[1,2])
=\mu_L(F\cap\coint{0,1})+\#(F\cap[1,2])
=\mu F$.}

\noindent {\bf (ii)} Jika $F\cap[1,2]$ tak berhingga, maka
$\phi^{-1}[F]\cap[1,2]$ juga demikian, sehingga

\Centerline{$\mu^*(\phi^{-1}[F])=\infty=\mu F$. \Qed}
}

Ini berarti bahwa jika kita menetapkan $J(x)=1$ untuk setiap
$x\in[0,2]$,
\cmmnt{

\Centerline{$\int_{\phi^{-1}[F]}J\,d\mu
=\mu_{\phi^{-1}[F]}(\phi^{-1}[F])
=\mu^*(\phi^{-1}[F])
=\mu F$}

\noindent untuk setiap $F\in\Sigma$, dan} $\phi$, $J$ memenuhi hipotesis
235A yang diubah. Namun, jika kita menetapkan $g=\chi\coint{0,1}$, maka
$g$ terintegralkan terhadap $\mu$, dengan $\int g\,d\mu=1$, sedangkan
\cmmnt{

\Centerline{$J(x)g(\phi(x))=1$ jika $x\in[0,1]\setminus A$, $0$
jika tidak,}

\noindent sehingga, karena $A\notin\Sigma$,} $J\times g\phi$
tidak\cmmnt{ terukur, dan karena itu (karena $\mu$ lengkap)
tidak} terintegralkan terhadap $\mu$.

\leader{235R}{Membalik \dvrocolon{beban}}\cmmnt{
Di sepanjang pembahasan di atas, saya telah menggunakan rumus

\Centerline{$\int J\times g\phi=\int g$,}

\noindent sebagai perluasan alami dari rumus

\Centerline{$\int g=\int g\phi\times \phi'$}

\noindent dalam kalkulus lanjut biasa. Namun kita juga dapat memindahkan
`turunan' $J$ ke sisi lain persamaan, sebagai berikut.

\medskip

\noindent}{\bf Teorema} Misalkan $(X,\Sigma,\mu)$, $(Y,\Tau,\nu)$
ruang-ruang ukur dan $\phi:X\to Y$, $J:Y\to\coint{0,\infty}$
fungsi-fungsi sedemikian sehingga $\int_FJ\,d\nu$ dan
$\mu\phi^{-1}[F]$ terdefinisi dalam $[0,\infty]$ dan sama untuk setiap
$F\in\Tau$. Maka $\int g\phi\,d\mu=\int J\times g\,d\nu$
setiap kali $g$ terukur secara virtual terhadap $\nu$ dan terdefinisi
$\nu$-hampir di mana-mana, serta salah satu integral
terdefinisi dalam $[-\infty,\infty]$.

\proof{ Misalkan $\nu_1$ ukuran integral tak tentu atas $\nu$ yang
didefinisikan oleh $J$, dan $\hat\mu$ pelengkapan $\mu$. Maka $\phi$
merupakan fungsi \imp\ untuk $\hat\mu$ dan $\nu_1$. \Prf\ Jika
$F\in\Tau$, maka $\nu_1F=\int_FJ\,d\nu=\mu\phi^{-1}[F]$; yakni,
$\phi$ merupakan fungsi \imp\ untuk $\mu$ dan $\nu_1\restr\Tau$. Karena
$\nu_1$ merupakan pelengkapan $\nu_1\restr\Tau$ (234Lb), $\phi$
merupakan fungsi \imp\ untuk $\mu$ dan $\nu_1$
(234Ba).\ \Qed

Tentu saja $\nu_1$ juga dapat kita pandang sebagai ukuran integral tak
tentu atas pelengkapan $\hat\nu$ dari $\nu$ (212Fb). Jadi, jika $g$
terukur secara virtual terhadap $\nu$ dan terdefinisi $\nu$-hampir di
mana-mana,

\Centerline{$\int J\times g\,d\nu
=\int J\times g\,d\hat\nu
=\int g\,d\nu_1
=\int g\phi\,d\hat\mu
=\int g\phi\,d\mu$}

\noindent jika salah satu dari kelima integral terdefinisi dalam
$[-\infty,\infty]$, menurut 235K, 235Gb, dan sekali lagi 212Fb.
}%end of proof of 235R

\exercises{
\leader{235X}{Latihan dasar (a)}
%\spheader 235Xa
Jelaskan apa yang disampaikan 235A ketika $X=Y$,
$\Tau=\Sigma$, $\phi$ merupakan fungsi identitas, dan
$\nu E=\alpha\mu E$ untuk setiap $E\in\Sigma$.
%235A

\spheader 235Xb Misalkan $(X,\Sigma,\mu)$ suatu ruang ukur, $J$ suatu
fungsi bernilai real nonnegatif yang terintegralkan pada $X$, dan
$\phi:D_{\phi}\to\Bbb R$ suatu fungsi terukur, dengan $D_{\phi}$
merupakan himpunan bagian koterabaikan dari $X$. Tetapkan

\Centerline{$g(a)=\int_{\{x:\phi(x)\le a\}}J$}

\noindent untuk $a\in\Bbb R$, dan misalkan $\mu_g$ ukuran
Lebesgue--Stieltjes yang terkait dengan $g$. Tunjukkan bahwa
$\int J\times f\phi\,d\mu=\int fd\mu_g$ untuk setiap fungsi bernilai real
yang terintegralkan terhadap $\mu_g$ dan dinotasikan dengan $f$.
%235A

\spheader 235Xc Misalkan $\Sigma$, $\Tau$, dan $\Lambda$
masing-masing merupakan aljabar-$\sigma$ dari himpunan-himpunan bagian $X$, $Y$, dan
$Z$. Kita katakan bahwa suatu fungsi $\phi:A\to Y$, dengan
$A\subseteq X$, terukur terhadap $(\Sigma,\Tau)$ jika
$\phi^{-1}[F]\in\Sigma_A$, yaitu aljabar-$\sigma$ subruang dari $A$,
untuk setiap $F\in\Tau$. Andaikan bahwa $A\subseteq X$,
$B\subseteq Y$, $\phi:A\to Y$ terukur terhadap $(\Sigma,\Tau)$ dan
$\psi:B\to Z$ terukur terhadap $(\Tau,\Lambda)$. Tunjukkan bahwa
$\psi\phi$ terukur terhadap $(\Sigma,\Lambda)$. Turunkan 235C.
%235C

\spheader 235Xd Misalkan $(X,\Sigma,\mu)$ suatu ruang
ukur dan $(Y,\Tau,\nu)$ suatu ruang ukur semihingga. Misalkan
$\phi:D_{\phi}\to Y$ dan $J:D_J\to\coint{0,\infty}$ merupakan
fungsi-fungsi yang didefinisikan pada himpunan-himpunan bagian
koterabaikan $D_{\phi}$, $D_J$ dari $X$ sedemikian sehingga
$\int J\times\chi(\phi^{-1}[F])d\mu$ ada $=\nu F$ setiap kali
$F\in\Tau$ dan $\nu F<\infty$. Misalkan $g$ suatu fungsi bernilai real
yang terukur terhadap $\Tau$, dan didefinisikan pada suatu himpunan
bagian koterabaikan dari $Y$. Tunjukkan bahwa $J\times g\phi$
terintegralkan terhadap $\mu$ jika dan hanya jika $g$ terintegralkan
terhadap $\nu$, dan kedua integralnya kemudian sama, asalkan
$(J\times g\phi)(x)$ kita tafsirkan sebagai
$0$ ketika $J(x)=0$ dan $g(\phi(x))$ tidak terdefinisi.
%235E

\spheader 235Xe Misalkan $(X,\Sigma,\mu)$ suatu ruang ukur dan
$E\in\Sigma$. Definisikan ukuran $\mu\LLcorner E$ pada $X$ dengan
menetapkan $(\mu\LLcorner E)(F)=\mu(E\cap F)$ setiap kali
$F\subseteq X$ sedemikian sehingga $F\cap E\in\Sigma$. Tunjukkan bahwa,
untuk setiap fungsi $f$ dari suatu himpunan bagian $X$ ke
$[-\infty,\infty]$, $\int fd(\mu\LLcorner E)=\int_Efd\mu$ jika salah
satunya terdefinisi dalam $[-\infty,\infty]$.
%235G

\sqheader 235Xf Misalkan $g:\Bbb R\to\Bbb R$ suatu fungsi
tak menurun yang kontinu mutlak pada setiap interval tertutup terbatas,
dan $\mu_g$ ukuran Lebesgue--Stieltjes terkait (114Xa, 225Xd). Tuliskan
$\mu$ untuk ukuran Lebesgue pada $\Bbb R$, dan misalkan
$f:\Bbb R\to\Bbb R$ suatu fungsi. Tunjukkan bahwa
$\int f\times g'\,d\mu=\int f\,d\mu_g$ dalam arti bahwa jika salah satu
integral ada, berhingga atau tak berhingga, maka yang lain juga ada,
dan keduanya kemudian sama.
%235J

\spheader 235Xg Misalkan $g:\Bbb R\to\Bbb R$ suatu fungsi
tak menurun dan $J$ suatu fungsi bernilai real nonnegatif yang
terintegralkan terhadap $\mu_g$, dengan $\mu_g$ ukuran
Lebesgue--Stieltjes yang didefinisikan dari $g$. Tetapkan
$h(x)=\int_{\ocint{-\infty,x}}J\,d\mu_g$ untuk setiap $x\in\Bbb R$,
dan misalkan $\mu_h$ ukuran Lebesgue--Stieltjes yang terkait dengan
$h$. Tunjukkan bahwa, untuk setiap $f:\Bbb R\to\Bbb R$,
$\int f\times J\,d\mu_g=\int fd\mu_h$, dalam arti bahwa jika salah satu
integral terdefinisi dalam $[-\infty,\infty]$ maka yang lain juga
demikian, dan keduanya kemudian sama.
%235J, 235Xf

\sqheader 235Xh Misalkan $X$ suatu himpunan dan $\lambda$, $\mu$,
$\nu$ tiga ukuran pada $X$ sedemikian sehingga $\mu$ merupakan ukuran
integral tak tentu atas $\lambda$, dengan turunan Radon-Nikod\'ym $f$,
dan $\nu$ merupakan ukuran integral tak tentu atas $\mu$, dengan
turunan Radon-Nikod\'ym $g$. Tunjukkan bahwa $\nu$ merupakan ukuran
integral tak tentu atas $\lambda$, dan bahwa $f\times g$ merupakan
turunan Radon-Nikod\'ym dari $\nu$ terhadap $\lambda$, asalkan
$(f\times g)(x)$ kita tafsirkan sebagai $0$ ketika $f(x)=0$ dan $g(x)$
tidak terdefinisi.
%235K

\spheader 235Xi Dalam 235M, jika $\nu$ tidak semihingga, tunjukkan
bahwa kita masih dapat menemukan $J$ sedemikian sehingga
$\int_{\phi^{-1}[F]}J\,d\mu=\nu F$ untuk setiap himpunan $F$ berukuran
berhingga. \Hint{gunakan `versi semihingga' dari $\nu$, sebagaimana
diuraikan dalam 213Xc.}
%235M

\spheader 235Xj Misalkan $(X,\Sigma,\mu)$ suatu ruang ukur
$\sigma$-hingga, dan $\Tau$ suatu subaljabar-$\sigma$ dari $\Sigma$.
Misalkan $\nu:\Tau\to\Bbb R$ suatu fungsional aditif terhitung
sedemikian sehingga $\nu F=0$ setiap kali $F\in\Tau$ dan $\mu F=0$.
Tunjukkan bahwa terdapat fungsi yang terintegralkan terhadap $\mu$, yaitu
$f$, sedemikian sehingga $\int_Ffd\mu=\nu F$ untuk setiap $F\in\Tau$.
\Hint{gunakan metode 235M, yang diterapkan pada bagian positif dan
negatif dari $\nu$.}
%235M

\spheader 235Xk Misalkan $(X,\Sigma,\mu)$ dan $(Y,\Tau,\nu)$
ruang-ruang ukur, dengan pelengkapan $(X,\hat\Sigma,\hat\mu)$ dan
$(Y,\hat\Tau,\hat\nu)$. Misalkan $\phi:D_{\phi}\to Y$,
$J:D_J\to \coint{0,\infty}$ merupakan fungsi-fungsi yang didefinisikan
pada himpunan-himpunan bagian koterabaikan dari $X$. Tunjukkan bahwa
jika $\int_{\phi^{-1}[F]}J\,d\mu=\nu F$ setiap kali $F\in\Tau$ dan
$\nu F<\infty$, maka $\int_{\phi^{-1}[F]}J\,d\mu=\nu F$ setiap kali
$F\in\hat\Tau$ dan $\hat\nu F<\infty$. Karena itu, atau dengan cara
lain, tunjukkan bahwa 235Pb berlaku untuk ruang-ruang tak lengkap
$(X,\Sigma,\mu)$ dan $(Y,\Tau,\nu)$.
%235P

\spheader 235Xl Misalkan $(X,\Sigma,\mu)$ suatu ruang ukur lengkap,
$Y$ suatu himpunan, $\phi:X\to Y$ suatu fungsi, dan
$\nu=\mu\phi^{-1}$ ukuran citra terkait pada $Y$. Misalkan $\nu_1$
suatu ukuran integral tak tentu atas $\nu$. Tunjukkan bahwa terdapat
ukuran integral tak tentu $\mu_1$ atas $\mu$ sedemikian sehingga
$\nu_1$ merupakan ukuran citra $\mu_1\phi^{-1}$.
%??

\spheader 235Xm
 Misalkan $(X,\Sigma,\mu)$ dan $(Y,\Tau,\nu)$ ruang-ruang ukur,
dan $\phi:X\to Y$ suatu fungsi \imp.
Misalkan $\nu_1$ suatu
ukuran integral tak tentu atas $\nu$. Tunjukkan bahwa terdapat
ukuran integral tak tentu $\mu_1$ atas $\mu$ sedemikian sehingga $\phi$
merupakan fungsi \imp\  untuk $\mu_1$ dan $\nu_1$.
%235G out of order query 

\spheader 235Xn\dvAnew{2013}
Misalkan $(X,\Sigma,\mu)$ dan $(Y,\Tau,\nu)$
ruang-ruang ukur, dan $\phi:X\to Y$ suatu fungsi \imp. Tunjukkan bahwa
$\overline{\int}h\phi\,d\mu\le\overline{\int}h\,d\nu$ untuk setiap
fungsi bernilai real $h$ yang didefinisikan hampir di mana-mana dalam
$Y$.
(Bandingkan 234Bf.)
%235G 

\leader{235Y}{Latihan lanjutan (a)}
%\spheader 235Ya
Tuliskan $\Tau$ untuk aljabar himpunan bagian Borel dari
$Y=[0,1]$, dan $\nu$ untuk pembatasan ukuran Lebesgue pada $\Tau$.
Misalkan $A\subseteq[0,1]$ suatu himpunan sedemikian sehingga $A$ dan
$[0,1]\setminus A$ keduanya memiliki ukuran luar Lebesgue $1$, dan
tetapkan $X=A\cup[1,2]$. Misalkan $\Sigma$ aljabar himpunan bagian Borel
relatif dari $X$, dan tetapkan $\mu E=\mu_A(A\cap E)$ untuk
$E\in\Sigma$, dengan $\mu_A$ ukuran subruang yang diinduksi pada $A$
oleh ukuran Lebesgue. Definisikan $\phi:X\to Y$ dengan menetapkan
$\phi(x)=x$ jika $x\in A$, dan $x-1$ jika $x\in X\setminus A$.
Tunjukkan bahwa $\nu$ merupakan ukuran citra $\mu\phi^{-1}$, tetapi,
dengan menetapkan $g=\chi([0,1]\setminus A)$, $g\phi$ terintegralkan
terhadap $\mu$ sedangkan $g$ tidak terintegralkan terhadap $\nu$.

\spheader 235Yb Misalkan $(X,\Sigma,\mu)$ suatu ruang probabilitas dan
$\Tau$ suatu subaljabar-$\sigma$ dari $\Sigma$. Misalkan $f$ suatu
fungsi nonnegatif yang terintegralkan terhadap $\mu$ dengan
$\int fd\mu=1$, sehingga ukuran integral tak tentunya $\nu$ merupakan
ukuran probabilitas. Misalkan $g$ suatu fungsi bernilai real yang
terintegralkan terhadap $\nu$ dan tetapkan $h=f\times g$, dengan
menafsirkan
$h(x)$ sebagai $0$ jika $f(x)=0$ dan $g(x)$ tidak terdefinisi.
Misalkan $f_1$, $h_1$ merupakan ekspektasi bersyarat dari $f$, $h$
pada $\Tau$ terhadap ukuran $\mu$, dan tetapkan $g_1=h_1/f_1$, dengan
menafsirkan $g_1(x)$ sebagai $0$ jika $h_1(x)=0$ dan $f_1(x)$ bernilai
$0$ atau tidak terdefinisi. Tunjukkan bahwa $g_1$ merupakan ekspektasi
bersyarat dari $g$ pada $\Tau$ terhadap ukuran $\nu$.
}%end of exercises

\endnotes{
\Notesheader{235} Saya melihat bahwa saya telah menggunakan banyak
ruang dalam bagian ini untuk perkara teknis; hipotesis teorema-teoremanya
berubah-ubah secara tak menentu, khususnya dengan kelengkapan yang
dipakai pada selang-selang yang tampaknya sembarang, dan gagasan-gagasan
berulang dalam pola yang tidak teratur. Tidak ada sesuatu yang mendalam,
dan sebagian besar pekerjaan terdiri atas verifikasi perincian dengan
susah payah. Masalah pada topik ini adalah kegunaannya. Hasil-hasil di
sini merupakan ungkapan abstrak integrasi dengan substitusi; penerapannya
terdapat di seluruh teori ukuran.
Karena itu, saya tidak dapat merasa cukup dengan teorema-teorema yang
mengungkapkan gagasan dasarnya secara elegan, tetapi harus mencari
rumusan-rumusan yang dapat saya kutip dalam argumen selanjutnya.

Saya berharap contoh-contoh dalam 235Bb, 235H, 235N, 235Q, 234Ya, dan
235Ya sedikit banyak akan meyakinkan Anda bahwa terdapat jebakan nyata
bagi orang yang lengah, dan bahwa verifikasi cermat yang dituliskan
begitu panjang memang diperlukan. Di sisi lain, untungnya dalam konteks
sederhana, ketika ukuran $\mu$, $\nu$ bersifat $\sigma$-hingga dan
transformasi $\phi$ merupakan isomorfisme Borel, tidak muncul kesulitan
yang tak teratasi, dan khususnya malapetaka ukuran citra tidak menyusahkan
kita. Namun untuk pembahasan lebih lanjut ke arah ini, saya merujuk Anda
pada penerapan dalam \S263, \S265, dan \S271, serta Volume 4.
}%end of notes

\frnewpage
